\documentclass[a4paper,12pt]{article}
\usepackage{iftex}
\ifPDFTeX
  \usepackage[T2A]{fontenc}
  \usepackage[utf8]{inputenc}
\else
  \usepackage{fontspec}
  \setmainfont{DejaVu Serif}
\fi
\usepackage[russian]{babel}
\usepackage{amsmath,amssymb,mathtools}
\usepackage{geometry}
\geometry{left=22mm,right=20mm,top=23mm,bottom=24mm}
\usepackage{microtype}
\usepackage{icomma}
\usepackage{booktabs,tabularx,array,longtable}
\usepackage{enumitem}
\usepackage{tikz}
\usetikzlibrary{arrows.meta,positioning,shapes.geometric,calc}
\usepackage{pgfplots}
\pgfplotsset{compat=1.18}
\usepackage{xcolor}
\usepackage{fancyhdr}
\usepackage{hyperref}
\usepackage{needspace}
\definecolor{accent}{HTML}{008080}
\definecolor{darkaccent}{HTML}{005A5A}
\definecolor{muted}{HTML}{5E6768}
\definecolor{linegray}{HTML}{D5DDDD}
\hypersetup{colorlinks=true,linkcolor=darkaccent,urlcolor=darkaccent,pdftitle={Чек-лист. Логарифмы и первые шаги в тригонометрии},pdfauthor={Лёвин Артём Александрович}}
\setlength{\parindent}{0pt}
\setlength{\parskip}{6pt}
\setlist[itemize]{topsep=3pt,itemsep=3pt,leftmargin=7mm}
\setlist[enumerate]{topsep=3pt,itemsep=4pt,leftmargin=8mm}
\renewcommand{\arraystretch}{1.25}
\pagestyle{fancy}
\fancyhf{}
\fancyhead[L]{\footnotesize\textcolor{muted}{Логарифмы и тригонометрия}}
\fancyhead[R]{\footnotesize\textcolor{muted}{Даниил Хрястов}}
\fancyfoot[C]{\footnotesize\thepage}
\renewcommand{\headrulewidth}{0.2pt}
\setlength{\headheight}{15pt}
\newcommand{\tang}{\mathop{\mathrm{tg}}\nolimits}
\newcommand{\cotang}{\mathop{\mathrm{ctg}}\nolimits}
\newcommand{\key}[1]{\textcolor{darkaccent}{\textbf{#1}}}
\newcommand{\task}[2]{\needspace{7\baselineskip}\medskip\textbf{#1. #2}\par}
\newcommand{\checkitem}[1]{\item[$\square$] #1}
\begin{document}
\thispagestyle{empty}
\begin{center}
\vspace*{23mm}
{\color{darkaccent}\Large МАТЕМАТИКА}\par
\vspace{19mm}
{\Huge Чек-лист}\par
\vspace{11mm}
{\color{accent}\LARGE Логарифмы: уравнения,\par графики и прикладные задачи}\par
\vspace{6mm}
{\color{accent}\Large Первые шаги в тригонометрии}\par
\vspace{18mm}
{\large Подготовка к ЕГЭ по профильной математике}\par
\vfill
{\large Даниил Хрястов}\par
\vspace{10mm}
{20 ноября 2026 г.}\par
\vspace{16mm}
{\small Лёвин Артём Александрович}\par
{\small эксклюзивно для Даниила Хрястова}\par
\vspace{14mm}
\begin{tikzpicture}
\draw[accent,line width=.85pt] (0,0) .. controls (1.1,.60) and (2.5,-.55) .. (3.8,0)
 .. controls (5.1,.60) and (6.4,-.55) .. (7.6,0);
\end{tikzpicture}
\end{center}
\clearpage
\section*{1. Что необходимо уметь}
\begin{itemize}[label=$\square$]
\checkitem{Читать логарифмическое выражение с правильными скобками и отличать основание от аргумента.}
\checkitem{Переходить между логарифмической и показательной формами, соблюдая ограничения.}
\checkitem{Решать уравнения с логарифмами, в том числе с разными основаниями.}
\checkitem{Аккуратно подставлять данные в прикладные формулы, сокращать коэффициенты и проверять ответ.}
\checkitem{Находить параметры логарифмической функции по двум точкам её графика.}
\checkitem{Восстанавливать таблицу значений тригонометрических функций для пяти основных углов.}
\end{itemize}
\section*{2. Логарифм: центральное определение}
Для $a>0$, $a\ne1$, $b>0$ запись
\[
\boxed{\log_a b=c\quad\Longleftrightarrow\quad a^c=b}
\]
означает: в какую степень нужно возвести число $a$, чтобы получить $b$? Это равносильный переход, который на занятии называли \emph{«глаз циклопа»}: основание логарифма возводится в степень, равную значению логарифма.
\[
\log_2 8=3\ \Longleftrightarrow\ 2^3=8,\qquad
\log_5 1=0\ \Longleftrightarrow\ 5^0=1.
\]
\key{Ограничения:} основание положительно и не равно единице, а аргумент строго положителен. Для $\log_a(x+2)$ необходимо $x>-2$.
\subsection*{Единица снаружи логарифма}
\[
1=\log_4 4,\qquad \log_4 5+1=\log_4 5+\log_4 4=\log_4 20.
\]
Здесь единица прибавляется \emph{ко всему логарифму}. Она не превращает число $5$ в число $6$.
\subsection*{Разные основания: внимательная проверка}
Разберём уравнение, обсуждавшееся при проверке домашней работы:
\[
\log_2 x=\log_4 5+1=\log_4 20=\frac12\log_2 20.
\]
Поэтому $\log_2 x=\log_2\sqrt{20}$, откуда
\[\boxed{x=\sqrt{20}=2\sqrt5}.\]
\key{Проверьте:} у логарифмов с разными основаниями нельзя просто приравнять аргументы. Сначала приведите основания к одному числу.
\section*{3. Логарифмические уравнения со скобками}
\subsection*{Пример. Изменение основания и ОДЗ}
\[
\log_2\!\left(\frac{x}{6}\right)=\log_{1/2}(x+1).
\]
ОДЗ: $x/6>0$ и $x+1>0$, то есть $x>0$. Так как $\log_{1/2}u=-\log_2u$, получаем
\[
\log_2\!\left(\frac{x}{6}\right)=-\log_2(x+1)
\ \Longleftrightarrow\ \log_2\!\left(\frac{x(x+1)}{6}\right)=0.
\]
Следовательно, $x(x+1)/6=1$, или $x^2+x-6=0$.
\[
D=1+24=25,\qquad x_{1,2}=\frac{-1\pm5}{2} \Rightarrow x_1=2,\ x_2=-3.
\]
Число $-3$ не удовлетворяет ОДЗ. \key{Ответ: $x=2$.}
\subsection*{Контроль знаков}
В формуле $x=(-b\pm\sqrt D)/(2a)$ минус стоит \emph{перед всем коэффициентом} $b$. Если $b=1$, то $-b=-1$. После вычисления корней обязательно подставьте допустимый результат в исходное уравнение.
\subsection*{Как избежать громоздкой арифметики}
\[
4200=42\cdot10^2,\qquad 0{,}7=7\cdot10^{-1},\qquad
0{,}4=4\cdot10^{-1}.
\]
В произведении степени десяти сокращаются: $10^{-1+2-1}=1$. Если числа крупные, разложите числитель и знаменатель на множители:
\[
\frac{10350}{575\cdot6}
=\frac{2\cdot3^2\cdot5^2\cdot23}{(5^2\cdot23)(2\cdot3)}=3.
\]
\key{Приём:} сокращайте общие множители до окончательного перемножения; всегда сохраняйте исходные числовые данные задачи.

\clearpage
\section*{4. Прикладные задачи: формула и единицы величин}
\subsection*{Тип 1. Температура по логарифмической модели}
Рассматривается расчётная модель
\[
X=\frac{\alpha c m}{\gamma}\log_2\!\left(\frac{T_{\mathrm{в}}-T_{\mathrm{р}}}{T-T_{\mathrm{р}}}\right),
\]
где $T_{\mathrm{р}}$ и $T_{\mathrm{в}}$ обозначают заданные температуры, $T$ --- искомую температуру, а остальные буквы --- известные параметры модели. Для положительного аргумента должно выполняться
\[
\frac{T_{\mathrm{в}}-T_{\mathrm{р}}}{T-T_{\mathrm{р}}}>0.
\]
При $T_{\mathrm{р}}=20$, $T_{\mathrm{в}}=100$, $\alpha=0{,}7$, $c=4200$, $m=0{,}4$, $\gamma=42$, $X=84$ имеем
\[
84=\frac{0{,}7\cdot4200\cdot0{,}4}{42}\log_2\frac{80}{T-20}
=28\log_2\frac{80}{T-20}.
\]
Значит,
\[
\log_2\frac{80}{T-20}=3\Longrightarrow\frac{80}{T-20}=8
\Longrightarrow T-20=10\Longrightarrow\boxed{T=30}.
\]
\key{Проверка:} $80/(30-20)=8>0$; $28\log_2 8=28\cdot3=84$.

\subsection*{Тип 2. Работа при сжатии газа}
Для рассматриваемой модели работы
\[
A=\alpha\nu T\log_2\!\left(\frac{V_1}{V_2}\right)
\]
на занятии использовались $A=10350$, $\alpha=5{,}75$, $\nu=2$, $T=300$, $V_1=40$.
\[
10350=5{,}75\cdot2\cdot300\log_2\frac{40}{V_2}
=3450\log_2\frac{40}{V_2}.
\]
Следовательно,
\[
\log_2\frac{40}{V_2}=3,\quad\frac{40}{V_2}=8,\quad
\boxed{V_2=5}.
\]
Поскольку $V_1$ и $V_2$ --- объёмы, здесь необходимы $V_1>0$, $V_2>0$. Значения подставляются в согласованных единицах; в ответе к текстовой задаче указывайте единицу искомой величины.

\clearpage
\section*{Алгоритм: решение прикладной задачи с логарифмом}
\thispagestyle{fancy}
\vspace{12mm}
\begin{center}
\tikzset{
startstop/.style={draw=accent,ellipse,align=center,minimum width=40mm,minimum height=10mm,fill=white},
process/.style={draw=linegray,rounded corners=2pt,align=center,text width=95mm,minimum height=13mm,fill=white},
decision/.style={draw=accent,diamond,aspect=2.8,align=center,inner sep=2pt,text width=43mm,fill=white},
arr/.style={-{Stealth[length=2.5mm]},draw=darkaccent,thick}}
\begin{tikzpicture}[node distance=9mm]
\node[startstop] (s) {Начало};
\node[process,below=of s] (a) {Выпишите данные, единицы величин и искомую переменную.};
\node[process,below=of a] (b) {Подставьте значения в модель; по возможности сократите множители.};
\node[decision,below=12mm of b] (c) {Аргумент логарифма положителен?};
\node[process,below=13mm of c] (d) {Выделите логарифм и перейдите к показательной записи.};
\node[process,below=of d] (e) {Решите полученное уравнение или пропорцию.};
\node[process,below=of e] (f) {Проверьте ОДЗ, исходную формулу, единицы и ответ.};
\node[startstop,below=of f] (g) {Ответ};
\node[process,right=20mm of c,text width=27mm] (no) {Проверьте запись и исходные данные.};
\draw[arr] (s)--(a);\draw[arr] (a)--(b);\draw[arr] (b)--(c);
\draw[arr] (c)--node[left,font=\small]{да}(d);
\draw[arr] (c.east)--node[above,font=\small]{нет}(no.west);
\draw[arr] (no.north) |- (b.east);
\draw[arr] (d)--(e);\draw[arr] (e)--(f);\draw[arr] (f)--(g);
\end{tikzpicture}
\end{center}
\clearpage
\section*{5. Логарифмическая функция по графику}
В задачах занятия встречались два разных вида функций:
\[
\boxed{f(x)=b+\log_a x}\qquad\text{и}\qquad
\boxed{g(x)=\log_a(x+b)}.
\]
В первом случае $b$ сдвигает график \emph{вверх/вниз}, во втором $b$ стоит \emph{внутри аргумента} и меняет область определения: для $f$ нужно $x>0$, для $g$ нужно $x+b>0$. Всегда $a>0$, $a\ne1$.
\subsection*{Пример 1. Вертикальный сдвиг}
Пусть график $f(x)=b+\log_a x$ проходит через точки $A(1;-3)$ и $B(4;-1)$. Найдём $f(32)$:
\[
\begin{cases}-3=b+\log_a1,\\-1=b+\log_a4.\end{cases}\qquad
\log_a1=0\Rightarrow b=-3.
\]
По второй точке $-1=-3+\log_a4$, то есть $\log_a4=2$, откуда $a^2=4$. Возможное отрицательное значение $a=-2$ не подходит, поскольку основание логарифма обязано быть положительным. Поэтому $a=2$.
\[
\boxed{f(x)=\log_2x-3},\qquad
\boxed{f(32)=5-3=2}.
\]
\begin{center}
\begin{tikzpicture}
\begin{axis}[width=.76\linewidth,height=64mm,xmin=0,xmax=7.2,ymin=-4.5,ymax=0.6,axis lines=middle,
 xlabel={$x$},ylabel={$y$},grid=major,grid style={draw=linegray},
 xtick={0,1,2,4,6},ytick={-4,-3,-2,-1,0},tick label style={font=\small},
 xlabel style={at={(axis description cs:1,.36)},anchor=west},ylabel style={at={(axis description cs:.08,1)},anchor=south},
 clip=false]
\addplot[darkaccent,thick,samples=200,domain=.09:7.1]{ln(x)/ln(2)-3};
\addplot[only marks,mark=*,mark size=2pt,accent] coordinates {(1,-3) (4,-1)};
\node[above right,font=\small] at (axis cs:1,-3) {$A(1;-3)$};
\node[above,font=\small] at (axis cs:4,-1) {$B(4;-1)$};
\end{axis}
\end{tikzpicture}
\end{center}
\subsection*{Пример 2. Сдвиг внутри логарифма}
Пусть $g(x)=\log_a(x+b)$ проходит через $A(-4;0)$ и $B(-1;2)$. Тогда
\[
\begin{cases}0=\log_a(b-4),\\2=\log_a(b-1).\end{cases}
\quad\Rightarrow\quad b-4=1\Rightarrow b=5.
\]
Подставляем $b=5$: $2=\log_a4$, поэтому $a^2=4$ и с учётом $a>0$ получаем $a=2$.
\[
\boxed{g(x)=\log_2(x+5)},\qquad \text{ОДЗ: }x>-5.
\]
\clearpage
\section*{6. Как восстановить функцию и проверить решение}
\begin{enumerate}
\item Считайте с графика две \emph{надёжные узловые точки} $(x_1;y_1)$ и $(x_2;y_2)$. Удобно выбирать точки с целыми координатами; близость к нулю сама по себе не гарантирует, что точки удобные.
\item Выпишите нужную формулу $f(x)$ и подставьте координаты как $f(x_i)=y_i$.
\item Используйте $\log_a1=0$ или $\log_a a=1$, если аргументы позволяют упростить систему.
\item Найдите $a$ и $b$, затем проверьте $a>0$, $a\ne1$ и область допустимых значений для обеих точек.
\item Подставьте требуемое значение $x$ только после восстановления всей функции.
\end{enumerate}
\subsection*{Если перед логарифмом есть множитель}
Рассмотрим функцию $h(x)=-2\log_a(x+2)-6$. Если $h(1)=-8$, то
\[
-2\log_a3-6=-8\Rightarrow -2\log_a3=-2\Rightarrow\log_a3=1
\Rightarrow a=3.
\]
Если теперь требуется найти $x$, при котором $h(x)=-12$, получим
\[
-2\log_3(x+2)-6=-12\ \Longrightarrow\ \log_3(x+2)=3
\ \Longrightarrow\ x+2=27\ \Longrightarrow\ \boxed{x=25}.
\]
Заметьте, что деление на отрицательное число меняет знак результата. Всегда проверяйте: $25>-2$, поэтому логарифм определён.
\subsection*{Памятка по внимательности}
\begin{itemize}
\item \textbf{Знаки:} перепроверьте переносы, $-b$ в формуле квадратного уравнения и деление на отрицательное число.
\item \textbf{Скобки:} различайте $b+\log_ax$ и $\log_a(x+b)$; множитель перед логарифмом относится ко всему выражению.
\item \textbf{Логарифмы:} не сравнивайте аргументы, пока основания не стали одинаковыми.
\item \textbf{ОДЗ:} проверяйте знаменатель, аргумент и основание; отбрасывайте посторонние корни.
\item \textbf{Математический язык:} система означает выполнение всех условий одновременно; совокупность --- выполнение хотя бы одного условия.
\end{itemize}
\section*{7. Тригонометрия: пять опорных углов}
\[
0\quad\frac{\pi}{6}\;(30^\circ)\quad\frac{\pi}{4}\;(45^\circ)\quad
\frac{\pi}{3}\;(60^\circ)\quad\frac{\pi}{2}\;(90^\circ).
\]
Для синуса по порядку запишите под корнем $0,1,2,3,4$ и разделите каждое значение на $2$; для косинуса используйте обратный порядок $4,3,2,1,0$.
\begin{center}\small
\begin{tabular}{@{}>{\raggedright\arraybackslash}p{26mm}*{5}{>{\centering\arraybackslash}p{22mm}}@{}}
\toprule
$\theta$&$0$&$\pi/6$&$\pi/4$&$\pi/3$&$\pi/2$\\\midrule
$\sin\theta$&$0$&$\frac12$&$\frac{\sqrt2}{2}$&$\frac{\sqrt3}{2}$&$1$\\[5pt]
$\cos\theta$&$1$&$\frac{\sqrt3}{2}$&$\frac{\sqrt2}{2}$&$\frac12$&$0$\\[5pt]
$\tang\theta$&$0$&$\frac{\sqrt3}{3}$&$1$&$\sqrt3$&---\\[5pt]
$\cotang\theta$&---&$\sqrt3$&$1$&$\frac{\sqrt3}{3}$&$0$\\\bottomrule
\end{tabular}
\end{center}
\[
\tang\theta=\frac{\sin\theta}{\cos\theta}\quad(\cos\theta\ne0),\qquad
\cotang\theta=\frac{\cos\theta}{\sin\theta}\quad(\sin\theta\ne0).
\]
\textbf{Почему $\tang(\pi/2)$ не определён?} Потому что $\cos(\pi/2)=0$, а на нуль делить нельзя. Аналогично не определён $\cotang0$.
\subsection*{Убрать корень из знаменателя}
\[
\tang\frac{\pi}{6}=\frac{\tfrac12}{\tfrac{\sqrt3}{2}}
=\frac1{\sqrt3}\cdot\frac{\sqrt3}{\sqrt3}=\boxed{\frac{\sqrt3}{3}}.
\]
\clearpage
\section*{8. Тренировочный блок — путь А}
\textcolor{muted}{10 задач от базового до повышенного уровня. В каждой задаче выписывайте ОДЗ или физические ограничения, если они существенны. Проводите отдельную проверку знаков и скобок.}
\subsection*{Базовый уровень}
\task{1}{Логарифм и отдельно прибавленная единица}
При проверке домашних решений выяснилось, что важно правильно понимать положение единицы в логарифмической записи. Решите уравнение
\[
\log_3 x=\log_9 7+1.
\]
Явно замените единицу логарифмом по основанию $9$, а затем приведите обе части к одному основанию. Укажите допустимые значения $x$ и точное значение корня.
\task{2}{Уравнение с противоположными логарифмами}
В учебном тренажёре дано уравнение
\[
\log_3\!\left(\frac{x}{12}\right)=\log_{1/3}(x+4).
\]
Найдите все действительные решения. До преобразований определите область допустимых значений; после решения квадратного уравнения проверьте оба найденных кандидата по ограничениям.
\subsection*{Средний уровень}
\task{3}{Температура технологической смеси}
В лабораторной расчётной модели величина $X$ зависит от температуры $T$ по формуле
\[
X=\frac{\alpha c m}{\gamma}\log_2\!\left(\frac{T_{\mathrm{в}}-T_{\mathrm{р}}}{T-T_{\mathrm{р}}}\right).
\]
Известно: температура отсчёта $T_{\mathrm{р}}=10^\circ\mathrm{C}$, верхняя температура $T_{\mathrm{в}}=90^\circ\mathrm{C}$, $\alpha=0{,}5$, $c=3600$, $m=0{,}6$, $\gamma=30$, а измеренное значение $X=72$. Найдите температуру $T$ в градусах Цельсия. Сначала сократите числовой коэффициент, затем выделите логарифм и проверьте найденную температуру подстановкой.
\task{4}{Объём газа после сжатия}
При сжатии газа положительные объёмы $V_1$ и $V_2$ связаны с работой $A$ учебной моделью
\[
A=\alpha\nu T\log_2\!\left(\frac{V_1}{V_2}\right).
\]
В опыте начальный объём $V_1=48$ литров, температура $T=250$, коэффициент $\alpha=4$, количество вещества $\nu=3$, а работа $A=9000$ (в согласованных единицах модели). Определите конечный объём $V_2$ в литрах. Проверьте, что ответ физически согласуется со сжатием.
\task{5}{Сокращение расчётного коэффициента}
При обработке результатов практической работы инженер получил дробь
\[
K=\frac{12480}{520\cdot6}.
\]
Сначала упростите её без калькулятора разложением числителя и знаменателя на простые множители; отдельно покажите, какие множители сократились. Затем вычислите $K$. Цель упражнения --- отработать вычислительный приём из прикладных задач.
\clearpage
\subsection*{Уровень выше среднего: графики логарифмов}
\task{6}{Функция с вертикальным сдвигом}
На координатной плоскости изображён график $f(x)=b+\log_a x$. По точкам с целочисленными координатами установлено, что график проходит через $P(1;-2)$ и $Q(9;0)$. Восстановите параметры $a$ и $b$, запишите формулу функции и найдите $f(81)$. Объясните, почему при нахождении основания допустима только положительная величина.
\task{7}{Горизонтальный сдвиг логарифмического графика}
По графику функции $g(x)=\log_a(x+b)$ удалось достоверно считать две точки: $P(-3;0)$ и $Q(1;2)$. Найдите основание $a$ и параметр $b$, определите область определения функции и вычислите $g(21)$. Помните, что параметр $b$ находится внутри аргумента, поэтому влияет на допустимые значения $x$.
\task{8}{Множитель перед логарифмом и обратный поиск аргумента}
Модель задаётся функцией $h(x)=-2\log_a(x+1)-4$ с неизвестным основанием $a$. Известно, что её график проходит через точку $P(2;-6)$. Восстановите основание логарифма. Затем определите, при каком значении $x$ функция принимает значение $-10$. При делении уравнения на $-2$ особенно внимательно следите за знаком правой части.
\subsection*{Первое повторение тригонометрии}
\task{9}{Проверка таблицы стандартных углов}
Перед решением экзаменационного варианта требуется восстановить значения синуса, косинуса и тангенса для стандартных углов. Без калькулятора вычислите
\[
S=2\sin\frac{\pi}{6}+\cos\frac{\pi}{3}-\tang\frac{\pi}{4}.
\]
Сначала запишите используемые значения из таблицы, затем выполните арифметические действия. Укажите точное значение $S$.
\task{10}{Тангенс и котангенс: точное значение}
При самопроверке таблицы стандартных углов сравнивают $\cotang(\pi/6)$ и утроенный $\tang(\pi/6)$. Вычислите без калькулятора
\[
R=\cotang\frac{\pi}{6}-3\tang\frac{\pi}{6}.
\]
Выведите оба тригонометрических значения через синус и косинус. Если появится дробь с корнем в знаменателе, избавьтесь от такого знаменателя стандартным домножением.
\clearpage
\section*{9. Ответы и самопроверка}
\small
\begin{longtable}{@{}p{11mm}p{41mm}p{100mm}@{}}
\toprule
№ & Ответ & Ключевая проверка \\ \midrule
\endhead
1 & $x=3\sqrt7$ & $\log_9 7+1=\log_9 63=\log_3\sqrt{63}$.\\
2 & $x=2$ & $x(x+4)=12$, корень $-6$ исключается по ОДЗ $x>0$.\\
3 & $T=30^\circ\mathrm C$ & Коэффициент $36$; $\log_2\bigl(80/(T-10)\bigr)=2$.\\
4 & $V_2=6$ л & $\log_2(48/V_2)=3$; объём уменьшился в $8$ раз.\\
5 & $K=4$ & $12480/(520\cdot6)=24/6=4$.\\
6 & $a=3$, $b=-2$, $f(81)=2$ & $\log_a9=2\Rightarrow a=3$ (с учётом $a>0$).\\
7 & $b=4$, $a=\sqrt5$, $g(21)=4$ & $\log_a5=2$, поэтому $a^2=5$; ОДЗ $x>-4$.\\
8 & $a=3$, $x=26$ & $\log_a3=1$; далее $\log_3(x+1)=3$.\\
9 & $S=\frac12$ & $2\cdot\tfrac12+\tfrac12-1=\tfrac12$.\\
10 & $R=0$ & $\cotang(\pi/6)=\sqrt3$ и $3\tang(\pi/6)=\sqrt3$.\\
\bottomrule
\end{longtable}
\normalsize
\section*{Итоговая памятка}
\begin{tabularx}{\linewidth}{@{}p{52mm}X@{}}
\toprule
\textbf{Если в задаче...} & \textbf{Ваш первый шаг}\\\midrule
Разные основания логарифмов & Приведите их к одному основанию.\\
Логарифм равен числу & Используйте $\log_a b=c\iff b=a^c$.\\
Сложная прикладная формула & Выпишите данные, приведите единицы, сократите коэффициенты.\\
Две точки графика & Подставьте обе пары координат, решите систему для $a,b$.\\
Логарифм со сдвигом & Проверьте, находится ли $b$ в аргументе или перед логарифмом.\\
Угол $0,\pi/6,\pi/4,\pi/3,\pi/2$ & Восстановите таблицу: для $\sin$ под корнем $0$--$4$, для $\cos$ --- обратно.\\
Готово решение & Отдельно проверьте знаки, скобки, ОДЗ и единицы.\\
\bottomrule
\end{tabularx}
\vfill
\begin{center}\footnotesize\textcolor{muted}{Лёвин Артём Александрович --- эксклюзивно для Даниила Хрястова}\end{center}
\end{document}
